\chapter{弱拓扑与 Banach--Alaoglu 定理}
\label{chap:weak-topologies}

范数收敛要求所有单位范数泛函检测到的误差同时变小；若只要求每个固定
泛函检测到的误差趋零，就得到较弱的收敛。损失的是一致性，可能获得的则是
紧致性。在使用这种交换之前，必须把拓扑定义清楚：一个空间的弱拓扑和它
作为某个对偶空间时的弱星拓扑，通常并不相同。

\section{对偶配对与初始拓扑}

设 \(X\) 为赋范空间。对每个 \(\phi\in X^*\) 有标量映射
\(x\mapsto\phi(x)\)。使这些映射全部连续的最弱拓扑称为弱拓扑，
记 \(\sigma(X,X^*)\)。它是前文初始拓扑构造的一个例子。
Hahn--Banach 定理保证 \(X^*\) 分离点，故此拓扑 Hausdorff。

在对偶空间上，只用评价映射 \(\phi\mapsto\phi(x)\)（\(x\in X\)）作检测，
得到弱星拓扑，记为 \(\sigma(X^*,X)\)。
而 \(X^*\) 自己的弱拓扑是 \(\sigma(X^*,X^{**})\)，通常使用更多检测泛函。
这里 \(X^{**}=(X^*)^*\)，不把它未经证明地认作 \(X\)。

\section{基本邻域}

弱拓扑中 \(x_0\) 的基本邻域为
\[
 \{x:|\phi_j(x-x_0)|<\varepsilon,\ j=1,\ldots,m\},
\]
其中 \(\phi_1,\ldots,\phi_m\in X^*\)，\(m\) 有限，\(\varepsilon>0\)。
弱星拓扑中 \(\phi_0\) 的基本邻域则为
\[
 \{\phi:|\phi(x_j)-\phi_0(x_j)|<\varepsilon,\ j=1,\ldots,m\}.
\]
有限个条件产生的邻域与无穷多个条件同时成立的范数球不可混同。
若 \(X\) 无限维，有限多个泛函的公共核含非零向量；基本弱邻域包含这个
方向的整条直线，因而范数无界。

这些邻域是凸的，且加法、数乘连续。因此弱拓扑和弱星拓扑都是局部凸向量
空间拓扑；“局部凸”在这里表示零点具有凸邻域基，不要求存在一个范数产生它。

\section{强、弱与弱星收敛}

记 \(x_\alpha\rightharpoonup x\)，若
\(\phi(x_\alpha)\to\phi(x)\) 对每个 \(\phi\in X^*\) 成立。
记 \(\phi_\alpha\overset{*}{\rightharpoonup}\phi\)，若
\(\phi_\alpha(x)\to\phi(x)\) 对每个 \(x\in X\) 成立。
指标可为序列，也可为一般有向集上的网。由有限基本邻域的描述，这恰好是
相应拓扑中的收敛。

范数收敛蕴含弱收敛；在 \(X^*\) 中，弱收敛又蕴含弱星收敛。
例如在 \(\ell^2\) 中，单位坐标向量 \(e_n\) 弱收敛到零，
因为每个对偶由 \(g\in\ell^2\) 表示，且 \(\langle e_n,g\rangle=g_n\to0\)；
但 \(\|e_n\|_2=1\)。在 \(\ell^1\) 中，同一坐标向量列不弱收敛到零，
因为常数列 \(\mathbf1\in\ell^\infty=(\ell^1)^*\) 的检测值恒为一。

\section{有界性与下半连续性}

\begin{proposition}\label{prop:weak-sequence-bounded}
赋范空间中的弱收敛序列范数有界。若 \(X\) Banach，则 \(X^*\) 中的
弱星收敛序列范数有界。
\end{proposition}
\begin{proof}
对弱收敛列 \(x_n\)，把它看成定义在 Banach 空间 \(X^*\) 上的评价泛函
\(Jx_n:\phi\mapsto\phi(x_n)\)。这些泛函逐点有界，一致有界原理与
\(\|Jx_n\|=\|x_n\|\) 给结论。第二条直接对定义在 \(X\) 上的泛函列应用
一致有界原理。
\end{proof}

此处写的是序列。收敛网的全部取值未必有界：收敛只控制最终部分，而有向集
中一个最终部分之外的指标未必只有有限个，不能照搬收敛数列的有界性证明。
例如将 \(\N\times\N\) 按坐标分别比较排序，令
\(u_{(0,j)}=j\)，而 \(u_{(i,j)}=0\)（\(i\ge1\)）。该实数网在指标
\((1,0)\) 之后恒为零，因而收敛到零，但其全部取值无界。

\begin{proposition}\label{prop:weak-lower-semicontinuity}
范数在弱拓扑下下半连续；凸集的范数闭包等于其弱闭包。
\end{proposition}
\begin{proof}
由对偶范数刻画，闭球为闭集的交
\(\{x:\|x\|\le r\}=\bigcap_{\phi\in B_{X^*}}\{x:|\phi(x)|\le r\}\)。
所以范数下半连续，尤其弱收敛序列满足
\(\|x\|\le\liminf_n\|x_n\|\)。弱闭集必范数闭，故范数闭包包含于弱闭包。
空凸集的两种闭包都是空集。对非空凸集的范数闭包以外的点，
严格分离定理给一个弱开半空间，将该点与集合
分开，于是它也不在弱闭包内。
\end{proof}

\section{对偶球的乘积表示}

每个 \(\phi\in B_{X^*}\) 都满足 \(|\phi(x)|\le\|x\|\)。
因此可将它嵌入紧圆盘（实情形为闭区间）的乘积
\[
 \Phi:B_{X^*}\longrightarrow\prod_{x\in X}D_x,\qquad
 \Phi(\phi)=(\phi(x))_{x\in X},\quad
 D_x=\{z\in\mathbb K:|z|\le\|x\|\}.
\]
这也可用极集语言表述：单位球 \(B_X\) 的极集为
\(B_X^\circledast=\{\phi:|\phi(x)|\le1\ (x\in B_X)\}=B_{X^*}\)。
这里专用符号 \({}^\circledast\) 与前章表示开球的 \({}^\circ\) 区分。
乘积坐标记录的是所有评价，而积拓扑仍只在每次邻域中约束有限个坐标。

\section{Banach--Alaoglu 定理}

\begin{theorem}\label{thm:banach-alaoglu}
任意赋范空间 \(X\) 的对偶闭单位球 \(B_{X^*}\) 在弱星拓扑下紧致。
\end{theorem}
\begin{proof}
由 \theoremref{thm:tychonoff-product}，上述乘积紧致。
其点 \(z=(z_x)\) 属于 \(\Phi(B_{X^*})\) 当且仅当满足全部线性等式
\[
 z_{x+y}=z_x+z_y,\qquad z_{\lambda x}=\lambda z_x
 \quad(x,y\in X,\ \lambda\in\mathbb K).
\]
每个等式定义乘积中的闭集，故共同满足它们的集合闭。反向确实成立：
这些等式给线性映射 \(x\mapsto z_x\)，而坐标圆盘约束给
\(|z_x|\le\|x\|\)，所以它是范数至多一的连续泛函。
最后，\(\Phi\) 的坐标映射正是弱星拓扑的定义映射，故它是到像集的
同胚。闭子集紧致，定理得证。
\end{proof}

该证明没有假设 \(X\) 完备，也没有使用反射性；所得是对偶球的弱星紧致。
如果要把“紧”改写成“每个序列可抽收敛子列”，还需下一节的度量化，
或另外的序列紧致性定理。

\section{可分情形}

\begin{theorem}\label{thm:weakstar-metrizable-ball}
若 \(X\) 可分，则 \(B_{X^*}\) 上的弱星拓扑可度量化。因而其中每个序列
都有弱星收敛子列。
\end{theorem}
\begin{proof}
可分度量空间的子空间可分，故 \(B_X\) 有至多可数稠密集。
允许重复，把它列为序列 \((x_n)_{n\ge1}\)；若 \(X=\{0\}\)，每项均取零。
定义
\[
 d(\phi,\psi)=\sum_{n\ge1}2^{-n}
       \min\{1,|\phi(x_n)-\psi(x_n)|\}.
\]
稠密性和泛函连续性保证正定性，三角不等式由截断标量距离得到。
给定 \(d\)-球的半径 \(\varepsilon>0\)，先取 \(N\) 使
\(\sum_{n>N}2^{-n}<\varepsilon/2\)。若前 \(N\) 个坐标误差均小于
\(\varepsilon/2\)，则 \(d(\phi,\psi)<\varepsilon\)。
这些有限条件给出弱星邻域，所以每个 \(d\)-球包含同中心的弱星邻域。

反向，若 \(0<\tau<1\) 且 \(d(\phi,\psi)<2^{-n}\tau\)，则
\(|\phi(x_n)-\psi(x_n)|<\tau\)。对 \(x\in B_X\)，选
\(\|x-x_n\|<\eta\)，有
\[
 |\phi(x)-\psi(x)|\le2\eta+|\phi(x_n)-\psi(x_n)|.
\]
给定 \(x\) 的评价误差限，先选 \(\eta\) 使第一项小于一半误差限，
再用上述 \(d\)-球控制第二项。一般的非零 \(x\in X\) 先归一化为
\(x/\|x\|\)，再把误差乘以 \(\|x\|\)；零向量不产生约束。
对有限个评价条件，取相应球半径的最小值，就得到包含于给定弱星邻域
的 \(d\)-球。这证明拓扑相同，而不只是两者序列收敛相同。
紧致度量空间序列紧致，得到子列结论。
\end{proof}

\section{紧致与序列紧致}

上述由稠密坐标控制一般评价的估计使用了单位球的统一范数界；
仅凭稠密坐标收敛不能在无界集合上重复该论证。这不排除有限维空间的
整个弱星拓扑可度量化。一般弱星拓扑可能需要不可数个
评价坐标，Tychonoff 定理给的是开覆盖紧致性，不能仅凭“它是一个积空间”
就使用对角抽取。
本书采用 ZFC，以上使用已建立的紧 Hausdorff 乘积定理，不在此声称这种
证明给出了最弱选择原则。下一章将证明 Banach 空间弱拓扑的
Eberlein--Šmulian 定理；它是额外的线性分析结果，并不适用于任意拓扑空间。

\section*{习题}
\addcontentsline{toc}{section}{习题}
\subsection*{A组}
\begin{enumerate}[label=\textbf{\thechapter.\arabic*.}]
\exerciseitem{weak-finite-dimensional}证明有限维空间中弱拓扑等于范数拓扑。
\exerciseitem{weak-neighborhood-unbounded}证明无限维赋范空间的每个弱零邻域都范数无界。
\exerciseitem{weak-unique-limit}证明弱极限唯一。
\exerciseitem{weakstar-coordinate-linfty}设 \((x_\alpha)\) 是 \(\ell^\infty=(\ell^1)^*\)
中的范数有界网。证明：它弱星收敛到 \(x\)，当且仅当每个坐标收敛到
\(x\) 的对应坐标；若只预设逐坐标极限存在，先证明这些极限组成
\(\ell^\infty\) 的元素。序列作为特例。
\end{enumerate}
\subsection*{B组}
\begin{enumerate}[label=\textbf{\thechapter.\arabic*.},resume]
\exerciseitem{weakstar-not-weak-linfty}令 \(s_n=(0,\ldots,0,1,1,\ldots)\)，前 \(n\) 项为零。证明它在 \(\ell^\infty\) 中弱星趋零，却不弱趋零。
\exerciseitem{weak-convex-mazur}若 \(x_n\rightharpoonup x\)，证明对每个 \(n\ge1\)，
尾集 \(\{x_k:k\ge n\}\) 的凸包中可选 \(y_n\) 使 \(\|y_n-x\|<1/n\)。
\exerciseitem{weak-operator-continuous}证明有界线性算子把弱收敛网送到弱收敛网。
\exerciseitem{weak-unit-sphere-closure}若 \(X\) 无限维，证明单位球面在弱拓扑下的闭包为闭单位球。
\end{enumerate}
\subsection*{C组}
\begin{enumerate}[label=\textbf{\thechapter.\arabic*.},resume]
\projectexerciseitem{weakstar-incomplete-unbounded}在 \(X=c_{00}\) 配上上确界范数时，令 \(\phi_n(x)=nx_n\)。
证明 \(\phi_n\overset{*}{\rightharpoonup}0\) 而 \(\|\phi_n\|\to\infty\)，解释与正文的有界性命题不矛盾。
\end{enumerate}
